% Run pdfLaTeX from the source directory. Keep the lookup for the activity's year.
% Exactly ONE worked example, ONE guided example, and FIVE independent exercises.
% Counts apply to the whole activity; never repeat these blocks per family.
% No instructional highlights, highlighted text/math, semantic highlight commands,
% or highlight legends anywhere, including optional review solutions.
% Strip source highlighting while preserving content and ordinary formatting.
% Avoid inherited highlight-producing macros (e.g. the shaded preamble boxed macro).
\makeatletter
\def\input@path{{../../../preamble/}{../../../../book2026-2027/preamble/}}
\makeatother
\input{preamble.tex}
\usepackage{needspace}
\usepackage{varwidth}

\geometry{letterpaper,margin=0.25in}
\pagestyle{empty}
\linespread{1.04}
\setlength{\parskip}{4pt}
\setlength{\tabcolsep}{6pt}
\setlength{\extrarowheight}{2pt}

\newif\ifprintsolutions
\printsolutionsfalse % Set \printsolutionstrue to include the review pages.

\newcolumntype{Y}{>{\centering\arraybackslash}X}
\newcommand{\stageheading}[1]{%
    \par\Needspace{6\baselineskip}\addvspace{0.65\baselineskip}%
    {\centering\large\bfseries #1\par}\nopagebreak[4]
}
\newcommand{\problemtitle}[3]{%
    \par\Needspace{4\baselineskip}\addvspace{0.4\baselineskip}%
    \noindent
    {\begin{varwidth}[t]{0.75\linewidth}%
         \bfseries\strut #1. #2%
     \end{varwidth}}%
    \hfill\textbf{#3}\par\nopagebreak[4]
}
\newcommand{\task}[2]{\par\noindent\textbf{#1} #2\par}
\newcommand{\responseline}{\par\noindent\rule{0pt}{1.55em}\hrulefill\par}
\newcommand{\workspace}[1]{\par\noindent\rule{0pt}{#1}\par}

\begin{document}

{\centering
    {\Large\bfseries [Criterion and activity title]\par}
    Global Economics \textbar{} [Grade] \textbar{} [Term] \textbar{} [Lesson]\par
}
\textbf{[Criterion code]} [Criterion statement.]

\stageheading{Equations and basic procedure}
[Notation, rules, ordered steps, and checks.]

% Exactly one fully solved worked example, with clear step-by-step reasoning.
\stageheading{Worked example}
\problemtitle{1}{[Worked problem title]}{Type A}
[Context, givens, units, and constraints.]
\par\noindent
\begin{minipage}[t]{0.48\linewidth}
    \task{Task 1.}{[Explain the method, substitution, and intermediate calculations.]}

    \begin{tabularx}{\linewidth}{@{}lY@{}}
        \toprule
        {[Data heading]} & [Result and units] \\
        \midrule
        {[Given value]} & [Calculated value] \\
        \bottomrule
    \end{tabularx}
\end{minipage}%
\hfill
\begin{minipage}[t]{0.48\linewidth}
    \task{Task 2.}{[Explain all remaining steps, verify the result, and give the final answer and interpretation in plain mathematics.]}
\end{minipage}\par

% Exactly one guided example; resize workspace for each actual task.
\stageheading{Guided Example}
\problemtitle{2}{[Guided problem title]}{Type A}
[Fresh context, givens, units, and constraints.]
\task{Task 1.}{[Structured prompt or starter calculation; leave key reasoning to students.]}
\workspace{4\baselineskip}
\task{Task 2.}{[Structured prompts for the remaining steps, check, and final interpretation.]}
\workspace{5\baselineskip}
\responseline

% Exactly five UNSOLVED independent exercises with appropriate workspace.
% Independent exercises 1--3 (problems 3--5): consolidate core skills.
% Independent exercises 4--5 (problems 6--7): slightly greater challenge,
% always within the same topic and assessment criterion.
\stageheading{Independent Practice}
\problemtitle{3}{[First independent problem title]}{Type A}
[Fresh context, givens, units, and constraints.]
\task{Task 1.}{[First task with fewer hints.]}
\task{Task 2.}{[Next task and interpretation.]}
\workspace{6\baselineskip}
\responseline

\problemtitle{4}{[Second independent problem title]}{Type A}
[Fresh context, givens, units, and constraints.]
\task{Task 1.}{[First task with fewer hints.]}
\task{Task 2.}{[Next task and interpretation.]}
\workspace{6\baselineskip}
\responseline

\problemtitle{5}{[Third independent problem title]}{[Family label]}
[Fresh context, givens, units, and constraints; consolidate core skills.]
\task{Task 1.}{[Core task with fewer hints; leave unsolved.]}
\task{Task 2.}{[Next task and interpretation; leave unsolved.]}
\workspace{6\baselineskip}
\responseline

\problemtitle{6}{[Fourth independent problem title]}{[Family label]}
[Fresh context, givens, units, and constraints; slightly more challenging within the same topic and criterion.]
\task{Task 1.}{[Modestly more demanding application of the core skills; leave unsolved.]}
\task{Task 2.}{[Next task and interpretation; leave unsolved.]}
\workspace{8\baselineskip}
\responseline

\problemtitle{7}{[Fifth independent problem title]}{[Family label]}
[Fresh context, givens, units, and constraints; slightly more challenging within the same topic and criterion.]
\task{Task 1.}{[Modestly more demanding application of the core skills; leave unsolved.]}
\task{Task 2.}{[Next task and interpretation; leave unsolved.]}
\workspace{8\baselineskip}
\responseline

% Optional review only: solutions to existing problems 2--7, no new examples.
% Student exercises above remain unsolved. No highlights on review pages either.
\ifprintsolutions
    \clearpage
    \stageheading{Solutions}
    \problemtitle{2}{[Guided problem title]}{Type A}
    \task{Task 1.}{[Show how to use the givens.]}
    \task{Task 2.}{[Show the reasoning, checks, and interpretation.]}

    \problemtitle{3}{[First independent problem title]}{Type A}
    \task{Task 1.}{[Show how to use the givens.]}
    \task{Task 2.}{[Show the reasoning, checks, and interpretation.]}

    \problemtitle{4}{[Second independent problem title]}{Type A}
    \task{Task 1.}{[Show how to use the givens.]}
    \task{Task 2.}{[Show the reasoning, checks, and interpretation.]}

    \problemtitle{5}{[Third independent problem title]}{[Family label]}
    \task{Task 1.}{[Show how to use the givens.]}
    \task{Task 2.}{[Show the reasoning, checks, and interpretation.]}

    \problemtitle{6}{[Fourth independent problem title]}{[Family label]}
    \task{Task 1.}{[Show how to use the givens.]}
    \task{Task 2.}{[Show the reasoning, checks, and interpretation.]}

    \problemtitle{7}{[Fifth independent problem title]}{[Family label]}
    \task{Task 1.}{[Show how to use the givens.]}
    \task{Task 2.}{[Show the reasoning, checks, and interpretation.]}
\fi

\end{document}

